% ID: SEL-D0077-OA-E01
% ESTADO: cotejada. Segunda pasada visual y borrador compilado revisados.
% Fuente: Selectividad/sel_2008_matematicas_aplicadas/Mat CCSS 5 - Examen y criterios de correción.pdf
% SHA-256: 6b6f35fe4872415542ee008d454a2502926c6ed0fffe87639ed13bed50759c6b
% Página: 1; opción A; ejercicio original 1.

De las restricciones que deben cumplir las variables $x$ e $y$ en un problema de programación lineal se deduce el siguiente conjunto de inecuaciones:
\[
2y-x\leq8,\quad x+y\geq13,\quad y+4x\leq49,\quad x\geq0,\quad y\geq0.
\]
\begin{enumerate}[label=\alph*)]
\item \textbf{(1.5 puntos)} Represente gráficamente el recinto determinado por estas inecuaciones.
\item \textbf{(1 punto)} Determine los vértices del recinto.
\item \textbf{(0.5 puntos)} Obtenga los valores extremos de la función $F(x,y)=3x-4y+12$ en ese recinto e indique en qué punto o puntos se alcanza cada extremo.
\end{enumerate}
